% Provenance (never printed): see sources/outline_sources.md and sources/book1/BRIEFS.md (g8/02).
\chapter{Dictionaries}\label{ch:g8:dictionaries}

Lina has two hundred contacts in her phone. She types ``Grandma'', and the
number appears at once. The phone does not read the \omterm{def:g7:lists:list}{list} from the top,
checking each name in turn: it \emph{looks the name up}. Python has a kind of
\omterm{def:g6:variables-and-input:value}{value} made for looking things up by name.

\begin{recall}
A \omterm{def:g7:lists:list}{list} keeps \omterm{def:g6:variables-and-input:value}{values} in order and is reached by \omterm{def:g7:lists:list}{index}, from 0
(\cref{ch:g7:lists}); finding an \omterm{def:g7:lists:list}{item} in it means walking through it. A
\omterm{def:g7:strings:string}{string} can be cut into words with \code{split} (\cref{ch:g7:strings}).
\end{recall}

\section{Looking a value up by its key}

\begin{definition}[Dictionary, key, entry]\label{def:g8:dictionaries:dictionary}\label{def:g8:dictionaries:key}\label{def:g8:dictionaries:entry}
A \emph{dictionary}\index{dictionary} is a \omterm{def:g6:variables-and-input:value}{value} that holds pairs: each pair,
an \emph{entry}\index{entry}, links a \emph{key}\index{dictionary key} to a
\omterm{def:g6:variables-and-input:value}{value}. The keys are all different; \code{d[key]} gives the \omterm{def:g6:variables-and-input:value}{value} linked to
\code{key}. A dictionary is written between braces, with a colon between
each key and its \omterm{def:g6:variables-and-input:value}{value}: \code{{"Sami": 12, "Noor": 13}}.
\end{definition}

\Needspace*{7\baselineskip}
\omcode{code/grade-8/02-dictionaries/listings/contacts.py}{1}{4}{A phone book as
a \omterm{def:g8:dictionaries:dictionary}{dictionary}.}\label{lst:g8:dictionaries:contacts}

\Needspace*{7\baselineskip}
\omoutput{code/grade-8/02-dictionaries/outputs/contacts.txt}{A \omterm{def:g6:variables-and-input:value}{value} found by its
\omterm{def:g8:dictionaries:dictionary}{key}; \code{in} tests whether a \omterm{def:g8:dictionaries:dictionary}{key} is
there.}\label{out:g8:dictionaries:contacts}

\begin{omfigure}
\begin{tikzpicture}[x=1cm, y=1cm, font=\footnotesize]
  \foreach \k/\v [count=\i from 0] in {Grandma/0612 345 678, Sami/0698 765 432, Noor/0655 501 020} {
    \node[draw=black!60, fill=ocBlue!12, minimum width=2.2cm, minimum height=0.6cm, font=\footnotesize\ttfamily] (k\i) at (0,-0.75*\i) {\k};
    \node[draw=black!60, fill=ocYellow!25, minimum width=2.8cm, minimum height=0.6cm, font=\footnotesize\ttfamily] (v\i) at (3.4,-0.75*\i) {\v};
    \draw[-{Stealth[length=2mm]}, black!50] (k\i) -- (v\i);
  }
  \node[font=\bfseries] at (0,0.65) {keys};
  \node[font=\bfseries] at (3.4,0.65) {values};
  \node[anchor=east, font=\footnotesize, ocRed] (q) at (-2.2,0) {\code{phone["Grandma"]}};
  \draw[-{Stealth[length=2mm]}, thick, ocRed] (q) -- (k0);
\end{tikzpicture}

{\small A \omterm{def:g8:dictionaries:dictionary}{dictionary} as two columns: each \omterm{def:g8:dictionaries:dictionary}{key} leads to its \omterm{def:g6:variables-and-input:value}{value}.}
\end{omfigure}

Asking for a \omterm{def:g8:dictionaries:dictionary}{key} that is not there is a \omterm{def:g7:finding-and-fixing-bugs:kinds}{runtime error}, a \code{KeyError}:

\Needspace*{5\baselineskip}
\omcode{code/grade-8/02-dictionaries/listings/key_error.py}{1}{2}{A name that is
not in the \omterm{def:g8:dictionaries:dictionary}{dictionary}.}\label{lst:g8:dictionaries:key-error}

\Needspace*{10\baselineskip}
\omoutput{code/grade-8/02-dictionaries/outputs/key_error_run.txt}{The last line
names the missing \omterm{def:g8:dictionaries:dictionary}{key}.}\label{out:g8:dictionaries:key-error}

Either test first with \code{in}, or use \code{get}, which returns a default
\omterm{def:g6:variables-and-input:value}{value} instead of stopping:

\Needspace*{6\baselineskip}
\omcode{code/grade-8/02-dictionaries/listings/get_default.py}{1}{3}{\code{get}
with a default.}\label{lst:g8:dictionaries:get}

\Needspace*{6\baselineskip}
\omoutput{code/grade-8/02-dictionaries/outputs/get_default.txt}{No error for
Leo.}\label{out:g8:dictionaries:get}

\begin{omfigure}
\begin{tikzpicture}[x=1cm, y=1cm, font=\scriptsize\ttfamily]
  \foreach \n [count=\i from 0] in {Ali,Bea,Cleo,Dan,Eli,Fay,Gus,Grandma} {
    \draw[black!60, fill=ocBlue!10] (1.3*\i,0) rectangle ++(1.3,0.55);
    \node at (1.3*\i+0.65,0.275) {\n};
  }
  \draw[-{Stealth[length=2mm]}, thick, ocRed, rounded corners] (0.3,0.75) -- (0.3,0.95) -- (9.75,0.95) -- (9.75,0.6);
  \node[font=\footnotesize\sffamily, ocRed, anchor=south] at (5.2,0.95) {a list: compare with every name in turn};
  \node[draw=black!60, fill=ocBlue!10, minimum width=1.6cm, minimum height=0.55cm] (g) at (4.6,-1.3) {Grandma};
  \node[font=\footnotesize\sffamily, anchor=east] (k) at (2.0,-1.3) {a dictionary};
  \draw[-{Stealth[length=2mm]}, thick, ocMint!60!black] (k) -- (g);
\end{tikzpicture}

{\small Finding Grandma: a walk through a \omterm{def:g7:lists:list}{list} against one step in a
\omterm{def:g8:dictionaries:dictionary}{dictionary}. How a \omterm{def:g8:dictionaries:dictionary}{dictionary} goes straight to its \omterm{def:g8:dictionaries:dictionary}{key} is explained in the
university books.}
\end{omfigure}

\section{Creating, changing, walking through}

An \omterm{def:g6:variables-and-input:assignment}{assignment} to \code{d[key]} adds an \omterm{def:g8:dictionaries:dictionary}{entry} if the \omterm{def:g8:dictionaries:dictionary}{key} is new, and changes
its \omterm{def:g6:variables-and-input:value}{value} if the \omterm{def:g8:dictionaries:dictionary}{key} is already there; \code{del d[key]} removes an \omterm{def:g8:dictionaries:dictionary}{entry}. A
\code{for} \omterm{def:g4:algorithms-repeat-and-decide:loop}{loop} over a \omterm{def:g8:dictionaries:dictionary}{dictionary} walks through its \omterm{def:g8:dictionaries:dictionary}{keys}, in the order in
which they were added. \code{d.items()} gives the \omterm{def:g8:dictionaries:dictionary}{key} and the \omterm{def:g6:variables-and-input:value}{value} together:
the \omterm{def:g4:algorithms-repeat-and-decide:loop}{loop} then has two \omterm{def:g6:variables-and-input:value}{variables}, separated by a comma.

\Needspace*{12\baselineskip}
\omcode{code/grade-8/02-dictionaries/listings/change_walk.py}{1}{9}{Adding,
changing, removing, walking.}\label{lst:g8:dictionaries:change-walk}

\Needspace*{9\baselineskip}
\omoutput{code/grade-8/02-dictionaries/outputs/change_walk.txt}{Three walks over
the same two entries.}\label{out:g8:dictionaries:change-walk}

% ledger: g8:python-dict-order
Python keeps a \omterm{def:g8:dictionaries:dictionary}{dictionary} in the order of its \omterm{def:g8:dictionaries:dictionary}{keys}' arrival; changing a \omterm{def:g6:variables-and-input:value}{value}
does not move its \omterm{def:g8:dictionaries:dictionary}{key}, which is why \code{apples} is still first.

\section{Counting with a dictionary}

A \omterm{def:g8:dictionaries:dictionary}{dictionary} of counts answers ``how many of each?''. For each \omterm{def:g7:lists:list}{item}, its
count is its old count plus one, and \code{get(item, 0)} gives 0 for an \omterm{def:g7:lists:list}{item}
not seen before:

\begin{center}
\code{counts[word] = counts.get(word, 0) + 1}
\end{center}

\Needspace*{10\baselineskip}
\omcode{code/grade-8/02-dictionaries/listings/count_letters.py}{1}{7}{Counting
letters.}\label{lst:g8:dictionaries:letters}

\Needspace*{6\baselineskip}
\omoutput{code/grade-8/02-dictionaries/outputs/count_letters.txt}{Each letter
once, with its count.}\label{out:g8:dictionaries:letters}

The same pattern counts words. This version prints the \omterm{def:g8:dictionaries:dictionary}{dictionary} after each
word, to show it grow:

\Needspace*{8\baselineskip}
\omcode{code/grade-8/02-dictionaries/listings/count_words.py}{1}{5}{Counting
words.}\label{lst:g8:dictionaries:words}

\omoutput{code/grade-8/02-dictionaries/outputs/count_words.txt}{A new \omterm{def:g8:dictionaries:dictionary}{key} for a
new word, a larger count for a word seen
before.}\label{out:g8:dictionaries:words}

\begin{omfigure}
\begin{tikzpicture}[x=1cm, y=1cm, font=\scriptsize\ttfamily]
  \foreach \w/\d [count=\i from 0] in {the/{the 1}, fox/{the 1, fox 1}, saw/{the 1, fox 1, saw 1}, the/{the 2, fox 1, saw 1}, hen/{the 2, fox 1, saw 1, hen 1}, and/{the 2, fox 1, saw 1, hen 1, and 1}} {
    \node[anchor=east, font=\footnotesize\ttfamily, ocRed!80!black] at (1.2,-0.5*\i) {\w};
    \node[anchor=west, draw=black!40, fill=ocYellow!15, inner sep=2pt] at (1.5,-0.5*\i) {\d};
  }
\end{tikzpicture}

{\small The counts after each of the first six words. Only the second
\code{the} changes a count instead of adding a \omterm{def:g8:dictionaries:dictionary}{key}.}
\end{omfigure}

\section{A record as a dictionary, a table as a list of records}

\begin{definition}[Record and field]\label{def:g8:dictionaries:record}\label{def:g8:dictionaries:field}
A \emph{record}\index{record} is a group of \omterm{def:g6:variables-and-input:value}{values} that describe one thing,
each under a name: a \emph{field}\index{field}. In Python, a record can be a
\omterm{def:g8:dictionaries:dictionary}{dictionary} whose \omterm{def:g8:dictionaries:dictionary}{keys} are the field names.
\end{definition}

A book of a school library has a title, an author and a year: three \omterm{def:g8:dictionaries:record}{fields}. A
\omterm{def:g7:lists:list}{list} of such \omterm{def:g8:dictionaries:record}{records} is a table, one \omterm{def:g8:dictionaries:record}{record} per row.

\omcode{code/grade-8/02-dictionaries/listings/books.py}{1}{13}{A small library:
a \omterm{def:g7:lists:list}{list} of three \omterm{def:g8:dictionaries:record}{records}.}\label{lst:g8:dictionaries:books}

\Needspace*{8\baselineskip}
\omoutput{code/grade-8/02-dictionaries/outputs/books.txt}{The \omterm{def:g8:dictionaries:record}{fields} of each
\omterm{def:g8:dictionaries:record}{record}, then a count.}\label{out:g8:dictionaries:books}

\begin{omfigure}
\begin{tikzpicture}[x=1cm, y=0.6cm, font=\footnotesize]
  \foreach \f/\x/\w in {title/0/3.6, author/3.6/2.8, year/6.4/1.2} {
    \draw[black!60, fill=ocInk!80] (\x,0) rectangle ++(\w,1);
    \node[white, font=\footnotesize\ttfamily] at (\x+\w/2,0.5) {\f};
  }
  \foreach \t/\a/\y [count=\i from 1] in {The Lighthouse Cat/Mara Quill/1998, Seven Bridges/Tom Alder/2011, Rain on Hilltop/June Okafor/2004} {
    \draw[black!60, fill=ocBlue!8] (0,-\i) rectangle ++(3.6,1); \node at (1.8,-\i+0.5) {\t};
    \draw[black!60, fill=ocBlue!8] (3.6,-\i) rectangle ++(2.8,1); \node at (5.0,-\i+0.5) {\a};
    \draw[black!60, fill=ocBlue!8] (6.4,-\i) rectangle ++(1.2,1); \node at (7.0,-\i+0.5) {\y};
    \pgfmathtruncatemacro{\j}{\i-1}
    \node[anchor=east, ocRed!80!black] at (-0.1,-\i+0.5) {\texttt{books[\j]}};
  }
\end{tikzpicture}

{\small The \omterm{def:g7:lists:list}{list} of \cref{lst:g8:dictionaries:books} as a table: the \omterm{def:g8:dictionaries:record}{field}
names head the columns, each \omterm{def:g8:dictionaries:record}{record} is a row.}
\end{omfigure}

\section{Exercises}

\begin{exercise}[$\star$]\label{exo:g8:dictionaries:1}
With \code{d = {"a": 1, "b": 2, "c": 3}}, what is \code{d["b"]}? What does
\code{d["d"]} do?
\end{exercise}

\begin{exercise}[$\star$]\label{exo:g8:dictionaries:2}
\code{d} is the empty \omterm{def:g8:dictionaries:dictionary}{dictionary} \code{{}}. After \code{d["x"] = 1} and then
\code{d["x"] = 5}, what are \code{len(d)} and \code{d["x"]}?
\end{exercise}

\begin{exercise}[$\star$]\label{exo:g8:dictionaries:3}
What does this \omterm{def:g4:programs-and-operating-system:program}{program} print?

\omcode{code/grade-8/02-dictionaries/exercises/walk_items.py}{1}{4}{A walk over
\code{items()}.}\label{lst:g8:dictionaries:walk-items}
\end{exercise}

\begin{exercise}[$\star$]\label{exo:g8:dictionaries:4}
Write a \omterm{def:g7:functions:function}{function} \code{capital(country)} that looks a country up in a
\omterm{def:g8:dictionaries:dictionary}{dictionary} of four capitals and returns \code{"not known"} for any other
country.
\end{exercise}

\begin{exercise}[$\star\star$]\label{exo:g8:dictionaries:5}
This \omterm{def:g4:programs-and-operating-system:program}{program} stops with \code{KeyError: 'mouse'}. Fix it so that a word
missing from the \omterm{def:g8:dictionaries:dictionary}{dictionary} is kept as it is.

\omcode{code/grade-8/02-dictionaries/exercises/translate_bug.py}{1}{6}{A word
for word translation.}\label{lst:g8:dictionaries:translate}
\end{exercise}

\begin{exercise}[$\star\star$]\label{exo:g8:dictionaries:6}
This \omterm{def:g4:programs-and-operating-system:program}{program} should count the words of \code{"a rose is a rose"}. It prints
\code{{'a': 1, 'rose': 1, 'is': 1}}. Why? Fix it.

\omcode{code/grade-8/02-dictionaries/exercises/counts_bug.py}{1}{4}{Counts that
never grow.}\label{lst:g8:dictionaries:counts-bug}
\end{exercise}

\begin{exercise}[$\star\star$]\label{exo:g8:dictionaries:7}
Write a \omterm{def:g7:functions:function}{function} \code{most_frequent(sentence)} that returns the word that
appears most often.
\end{exercise}

\begin{exercise}[$\star\star$]\label{exo:g8:dictionaries:8}
Write a \omterm{def:g7:functions:function}{function} \code{invert(d)} that swaps the \omterm{def:g8:dictionaries:dictionary}{keys} and the \omterm{def:g6:variables-and-input:value}{values} of a
\omterm{def:g8:dictionaries:dictionary}{dictionary} whose \omterm{def:g6:variables-and-input:value}{values} are all different. Why must they be different?
\end{exercise}

\begin{exercise}[$\star\star$]\label{exo:g8:dictionaries:9}
Without running it, how many entries does a word-count \omterm{def:g8:dictionaries:dictionary}{dictionary} have for
\code{"a rose is a rose is a rose"}? What are the counts?
\end{exercise}

\begin{exercise}[$\star\star\star$]\label{exo:g8:dictionaries:10}
Write a \omterm{def:g7:functions:function}{function} \code{oldest(books)} that returns the \omterm{def:g8:dictionaries:record}{record} of the oldest
book of a \omterm{def:g7:lists:list}{list} like the one of \cref{lst:g8:dictionaries:books}.
\end{exercise}

\begin{exercise}[$\star\star\star$]\label{exo:g8:dictionaries:11}
Votes for $c$ candidates are written on $b$ ballots. One method counts each
candidate's votes by walking all the ballots; the other walks the ballots once
with a \omterm{def:g8:dictionaries:dictionary}{dictionary} of counts. How many names does each method look at? Which
is better for 1\,000 ballots and 50 candidates, and by how much?
\end{exercise}

\begin{exercise}[$\star\star\star$]\label{exo:g8:dictionaries:12}
Write a \omterm{def:g7:functions:function}{function} \code{group_by_first_letter(words)} that returns a \omterm{def:g8:dictionaries:dictionary}{dictionary}
whose \omterm{def:g8:dictionaries:dictionary}{keys} are first letters and whose \omterm{def:g6:variables-and-input:value}{values} are \omterm{def:g7:lists:list}{lists} of the words that
start with each.
\end{exercise}

\section{Problem: Counting the Votes}

\begin{problem}[{Weekend problem --- two ways to count an election}]\label{pb:g8:dictionaries:1}
A class of 40 pupils elects its delegate among six candidates. Each pupil
writes one name on a ballot, and the ballots are kept in a \omterm{def:g7:lists:list}{list}.

\omcode{code/grade-8/02-dictionaries/listings/ballots.py}{1}{17}{The ballots of
the class, and a rule to make up a larger
election.}\label{lst:g8:dictionaries:ballots}

\needspace{4\baselineskip}\medskip\noindent\textbf{Part I --- By hand.}
\begin{enumerate}
  \item Tally the first 12 ballots by hand. Who leads?
  \item How did you keep the counts? What would you do with a name not yet on
    your tally?
\end{enumerate}

\needspace{4\baselineskip}\medskip\noindent\textbf{Part II --- Two programs.} The first \omterm{def:g4:programs-and-operating-system:program}{program} counts
each candidate's votes with a walk through all the ballots. A \omterm{def:g7:functions:function}{function} may
return two \omterm{def:g6:variables-and-input:value}{values} separated by a comma; the caller receives them the same
way, \code{results, looked_at = ...}.

\omcode{code/grade-8/02-dictionaries/listings/votes_list.py}{1}{22}{One walk per
candidate.}\label{lst:g8:dictionaries:votes-list}

\begin{enumerate}[resume]
  \item What does \code{looked_at} count? How often does the inner \omterm{def:g4:algorithms-repeat-and-decide:loop}{loop} run?
  \item Run it. Who wins, and how many names were looked at?
\end{enumerate}

The second \omterm{def:g4:programs-and-operating-system:program}{program} walks the ballots once:

\omcode{code/grade-8/02-dictionaries/listings/votes_dict.py}{1}{31}{One walk,
with a \omterm{def:g8:dictionaries:dictionary}{dictionary}.}\label{lst:g8:dictionaries:votes-dict}

\begin{enumerate}[resume]
  \item Which line is the counting pattern of this chapter?
  \item How does \code{winner_and_margin} find the second best?
  \item Run it. Do both \omterm{def:g4:programs-and-operating-system:program}{programs} agree? How many names were looked at now?
  \item Why does the \omterm{def:g8:dictionaries:dictionary}{dictionary} \omterm{def:g7:lists:list}{list} the candidates in a different order from
    \code{CANDIDATES}?
\end{enumerate}

\needspace{4\baselineskip}\medskip\noindent\textbf{Part III --- A larger election.} A third \omterm{def:g4:programs-and-operating-system:program}{program} runs
both methods on the class, then on a made-up election of 400 ballots for 20
candidates:

\omcode{code/grade-8/02-dictionaries/listings/votes_compare.py}{1}{17}{Both
methods on two elections.}\label{lst:g8:dictionaries:compare}

\begin{enumerate}[resume]
  \item Predict the four numbers before running it.
  \item Run it. Write a rule for each method in terms of the number of
    ballots $b$ and of candidates $c$.
  \item A country votes with 40 million ballots for 12 \omterm{def:g7:lists:list}{lists}. How many names
    would each method look at?
  \item State the result.
\end{enumerate}
\end{problem}
